% `verbatim-macros.tex`
%
% These are macros to make code extracts work, which is much trickier than it
% looks. There are LaTeX packages for this, but often those have tricky
% restrictions (can't be used in `\centerline`, `\footnote`, etc.). So we'll
% write our own macros, and cope with their issues by compiling around those
% issues (see the Inweb source code). In the end, this is code going back to
% Knuth's handling of verbatim code in the TeXBook, though slightly tweaked.
%
\newskip\ttglue
%
\chardef\other=12
\def\codetext{\begingroup\codetextspacing\frenchspacing
  \catcode`\\=\other
  \catcode`\{=\other
  \catcode`\}=\other
  \catcode`\$=\other
  \catcode`\&=\other
  \catcode`\#=\other
  \catcode`\%=\other
  \catcode`\~=\other
  \catcode`\_=\other
  \catcode`\^=\other
  \obeyspaces \obeylines \tt}
%
\def\|{\leavevmode\hbox{\tt\char`\|}} % vertical line
\outer\def\begintt{$$\let\par=\endgraf \codetext \parskip=0pt
  \catcode`\|=0 \rightskip-5pc \ttfinish}
{\catcode`\|=0 |catcode`|\=\other % | is temporary escape character
  |obeylines % end of line is active
  |gdef|ttfinish#1^^M#2\endtt{#1|vbox{#2}|endgroup|bodytext$$}}
\catcode`\|=\active
%
{\obeylines \gdef|{\codetext \spaceskip\ttglue \let^^M=\  \let|=\endgroup}}
%
\catcode`@=11
\def\beginlines{\par\begingroup\nobreak\medskip\parindent=0pt \obeylines\nobreak \everypar{\strut}}
\def\beginlinestighter{\par\begingroup\nobreak\parindent=0pt \kern1pt\nobreak \obeylines \everypar{\strut}}
\def\endlines{\kern1pt\endgroup\medbreak\noindent}
\catcode`@=12
%